\documentclass[12pt,aps,floats,showpacs,amssymb,tightenlines]{revtex4}



\usepackage{amsmath}
\usepackage{amsfonts}
\usepackage{amscd}
\usepackage{epsfig}
\usepackage{amssymb}
\usepackage{tabularx}

\renewcommand{\a}{\alpha}
\renewcommand{\b}{\beta}
\newcommand{\g}{\gamma}
\newcommand{\e}{\epsilon}
\newcommand{\m}{\mu}
\renewcommand{\d}{\delta}
\newcommand{\w}{\omega}
\renewcommand{\t}{\tau}
\newcommand{\p}{\wedge}
\newcommand{\A}{\bar{A}}
\newcommand{\R}{{\cal{R}}}
\newcommand{\s}{ Schwarzschild }


\begin{document}

\title{Instability of the negative mass Schwarzschild naked singularity}
\author{Reinaldo J. Gleiser and Gustavo Dotti}
\affiliation{Facultad de Matem\'atica, Astronom\'{\i}a y F\'{\i}sica,
Universidad Nacional de C\'ordoba, Ciudad Universitaria,
(5000) C\'ordoba, Argentina}
\email{gdotti@famaf.unc.edu.ar}


\begin{abstract}
We study the negative mass  Schwarzschild spacetime, which has a
naked singularity, and show that it is perturbatively unstable. This
is achieved by first introducing a modification of the well known Regge - Wheeler - Zerilli approach to black hole perturbations to allow for the presence of a ``kinematic'' singularity that arises for negative masses, and then exhibiting exact exponentially growing  solutions to the
linearized Einstein's equations. The perturbations are smooth everywhere and behave nicely around the
singularity and at infinity. In particular, the first order variation of the
scalar invariants can be made everywhere arbitrarily small as compared to the zeroth order terms.
Our approach is also compared to a recent analysis that leads to a different conclusion regarding the stability of the negative mass  Schwarzschild spacetime. We also comment on the relevance of our results to the stability of more general negative mass, nakedly singular spacetimes.
\end{abstract}
\pacs{04.50.+h,04.20.-q,04.70.-s, 04.30.-w}


\maketitle




\noindent
\section{Introduction: the negative mass Schwarzschild solution}

The positive mass theorems of general relativity establish that, under plausible assumptions,
asymptotically flat solutions of the Einstein equations with physically acceptable matter sources
cannot have negative total mass. Another well known feature of general relativity, compatible with
the positivity theorems, is the possibility of the presence, or eventual formation, of singularities
 in the spacetime manifold. However, in the most relevant examples where this happens, the singularity
 is surrounded by a horizon, so that it is unobservable to an observer outside the horizon. That this
 should happen generically in all physically relevant situations is the essential content of the versions
 of the Cosmic Censorship hypothesis. On the other hand, from a strictly mathematical point of view,
 there is plenty of solutions of the Einstein equations that are asymptotically flat and present a
  negative total mass, but, in general, they are also nakedly singular, which would make them physically
   unacceptable. A simple and particularly interesting example of this type of solutions is the
   negative mass Schwarzschild spacetime. Nevertheless, for different reasons, these solutions
   are not expected to arise from physically acceptable regular initial data. One should notice,
   however, that what is physically acceptable may change as more data is available. What we
   have in mind here is the recent suggestion of the existence of ``phantom matter", that is some
   field that carries {\em negative kinetic energy}, as a possible explanation for the observed
    acceleration of the scale factor of the Universe \cite{pm}.

A different, but equally important consideration regarding negative masses is the manner
in which they interact with positive masses. This gives rise to different types of
instabilities that have been analyzed in the literature. We refer to \cite{ghi} for a
recent review. It is rather remarkable, however, that until recently \cite{ghi}, the stability
of the negative mass Schwarzschild solution, that, in a loose sense, may be considered as representing
the field of an isolated negative mass source, had not been considered in the literature, although the
 essential machinery for this analysis, at least at the linear level, had been available for some time.
 The analysis carried out in \cite{ghi} leads to the conclusion that if certain boundary conditions,
 assumed suitable for the problem, are imposed, the negative mass Schwarzschild singularity should be considered as
 stable under linear perturbations. However, as discussed in Section \ref{ns}, the resulting solutions of Zerilli's equation, that are used to prove stability in
 \cite{ghi}, are only twice differentiable functions  of the radial coordinate. As a consequence, the
 perturbed metrics fail to be solutions of the vacuum of Einstein's equations, and they
 actually satisfy the field equations of a singular distribution of sources, unrelated to the negative mass Schwarzschild singularity.
  Given this situation, we undertake in this paper
 a new analysis of the linear stability of the negative mass Schwarzschild solution. This is based on
 a modification of the standard Regge-Wheeler-Zerilli formalism (reviewed in Section \ref{rw})
  that simplifies and makes more
 transparent the analysis of the perturbative equations for negative mass. This is given in Section
 \ref{rwzg}.
  Next, we consider the existence of ``zero modes'', i.e., static solutions of the perturbation equations.
   We give the general solution, that for each mode depends on two integration constants, and prove that
    there is no choice of these constants such that the solution can be considered as a ``small"
    perturbation. We give also a physical interpretation for these solutions that supports our
    conclusions. In Section \ref{us}, based on this modified form of the formalism, and this is the
     main result presented in this paper, we prove that under quite general boundary conditions,
      {\em all} scalar perturbative modes are unstable. Concurrently in the Appendix, we prove
      that our main perturbation equation, that replaces the Sch\"odinger like equation for the
      standard Zerilli function, has the remarkable property that {\em all} its solution are
      regular at the ``kinematic'' singularity that appears in the Regge-Wheeler-Zerilli formalism
       for negative mass. Our results and possible consequences are finally summarized in Section \ref{cc}.





\section{Perturbation equations for the spherically symmetric Schwarzschild solution.} \label{rw}

The vacuum solution of Einstein's equations for the exterior of a bounded spherically symmetric source,
first obtained by Schwarzschild, may be written in the form,
\begin{equation}
\label{Schw1}
ds^2 =  -\left(1-\frac{2 M}{r}\right) dt^2 +
 \left(1-\frac{2 M}{r}\right)^{-1} dr^2   +r^2 \left( d \theta^2 +  \sin^2(\theta)
d \phi^2 \right)
\end{equation}
where $\theta$ and $\phi$ are standard coordinates on the unit sphere, and $M$ is a real
(positive or negative) constant, that can be identified with the Newtonian mass of the source.
The range of $r$ is limited by the singularities of the metric coefficients in (\ref{Schw1}).
 For $M>0$ the singularity at $r=2M$ corresponds to a regular horizon, and the solution may be
extended to $r<2M$. The resulting space time is known as a ``black
hole''. The horizon hides
  a curvature singularity for $r=0$, rendering the portion of space time outside the horizon
  globally hyperbolic. Similarly, for negative mass ($M<0$) there is a singularity for $r=0$,
   but in this case there is no horizon, the singularity is ``naked'', and the resulting space
    time is not gobally hyperbolic.

The rather strange properties of the spacetime described by
(\ref{Schw1}) led to question the possibility of their physical
realization. Since the solution (at least for $r>2M$) is static, a
first step in this direction, carried out by Regge and Wheeler \cite{RW}, and
completed later by Zerilli \cite{Zerilli}, was aimed at establishing their
stability, through an analysis of the behavior of linear
perturbations. The Einstein equations lead in this case to a set of
coupled linear partial differential equations for the functions
representing the perturbations. Because of the spherical symmetry of
(\ref{Schw1}) this functions may be grouped into ``even'' (or
scalar) and ``odd'' (or vector) sets that do not couple to each
other, and are characterized by their dependence on the angular
coordinates $(\theta,\phi)$. Still, the resulting perturbative
equations contain, in general, some redundant information that can
be eliminated through a choice of ``gauge''.
A particular choice, that keeps full generality, was selected in
\cite{RW}, \cite{Zerilli} and is known as the ''Regge-Wheeler gauge''. Independently of the gauge issue, the stability of the Schwarzschild space time, well known for $M>0$, was also established for $M<0$ in \cite{ghi}. For this reason we shall restrict our analysis to the scalar
perturbations.  In the Regge-Wheeler gauge, the scalar (even)
perturbations, for each angular mode $(\ell,m)$,  are fully
described by  four functions $H_0(r,t)$, $H_1(r,t)$, $H_2(r,t)$ and
$K(r,t)$, in terms of which the perturbed metric takes the form,
\begin{eqnarray}
\label{RW1}
ds^2 & = & -\left(1-\frac{2M}{r}\right)\left(1- \e H_0 Y_{\ell,m}\right)dt^2 +
2 \: \e \;  H_1 Y_{\ell,m} dt dr+\left(1-\frac{2M}{r}\right)^{-1}\left(1+ \e H_2 Y_{\ell,m}\right)dr^2
\nonumber \\
& & +r^2\left(1+ \e K\, Y_{\ell,m}\right) \left( d \theta^2 +  \sin^2(\theta)
d \phi^2 \right)
\end{eqnarray}
where $\epsilon$ is an auxiliary parameter, useful in keeping track of the order of the perturbation.
$Y_{\ell,m}=  Y_{\ell,m}(\theta,\phi)$ are standard spherical
harmonics on the sphere, and, on account of the spherical symmetry, we may restrict to $m=0$ without loss of generality. The linearized Einstein equations for the
metric (\ref{RW1}) imply $H_0(r,t)=H_2(r,t)$, and a set of coupled
differential equations for $H_1$, $H_2$ and $K$. This can be solved
in terms of the Zerilli function $\psi_Z(r,t)$, by the replacements,
\begin{eqnarray}
\label{RW2}
K & = & f(r) \psi_Z+\left(1-\frac{2M}{r}\right)\frac{\partial \psi_Z}{\partial r} \nonumber \\
H_1 & = & h(r)\frac{\partial \psi_Z}{\partial t}+r\frac{\partial^2 \psi_Z}{\partial t \partial r}
 \nonumber \\
H_2 & = & \frac{\partial }{\partial r}\left[\left(1-\frac{2M}{r}\right)\left(h(r) \psi_Z
+r\frac{\partial \psi_Z}{ \partial r} \right)\right] -K
\end{eqnarray}
where,
\begin{eqnarray}
\label{RW2a}
f(r) & = &  \frac{\lambda(\lambda+1)r^2+3 \lambda M r + 6 M^2}{r^2(\lambda r + 3 M)}\nonumber \\
h(r) & = &  \frac{\lambda r^2 - 3 \lambda r M -3 M^2}{(r-2M)(\lambda r + 3M)}\nonumber \\
\lambda & = &  \frac{(\ell-1)(\ell+2)}{2}
\end{eqnarray}
provided $\psi_Z$ is a solution of the Zerilli equation,
\begin{equation}
\label{RW3}
\frac{\partial^2 \psi_Z}{\partial t^2}-\frac{\partial^2 \psi_Z}{\partial r_*^2} +V(r) \psi_z = 0
\end{equation}
where,
\begin{equation}
\label{poten1}
V(r) = 2 \left(1 -\frac{2 M}{ r}\right) {\frac{ \lambda^2 r^2
\left[(\lambda+1) r + 3 M \right] + 9 M^2 (\lambda r +M)}{ r^3
(\lambda r + 3 M)^2} } ,
\end{equation}
and $r_*$ is the ``tortoise'' coordinate, related to $r$ by,
\begin{equation}
\label{RW4}
\frac{dr_*}{dr}=\left(1-\frac{2M}{r}\right)^{-1}
\end{equation}

In what follows, for simplicity, we choose the integration constant relating $r$ and $r_*$ such that,
\begin{equation}
\label{RW4a}
r_* = r +2 M \left[\ln (r-2M) -\ln(2|M|)\right]
\end{equation}
Therefore, for $M<0$, we have $ r_*=0$ for $r=0$.

We notice that the relations (\ref{RW2}) can be inverted, giving,
\begin{equation}
\label{RW3a}
\psi_Z(r,t)  = \frac {r (r-2M)}{ ( \lambda+1 )  ( \lambda
\,r+3\,M ) } \left(H_2
 - r  \frac {\partial K}{\partial r}
\right)
    +\frac{r}{\lambda+1} K
\end{equation}
Therefore, there is a one to one correspondence between $\psi_Z$ and
the Regge-Wheeler gauge functions, $H_0$, $H_1$, $H_2$ and $K$.\\

For the stability analysis one considers a time dependence of the form,
\begin{equation}
\label{psitopsi1}
\psi_Z(t,r)=e^{i \omega t} \psi(r)
\end{equation}
which turns (\ref{RW3}) into the Schr\"odinger-like form,
\begin{equation}
\label{RW3b}
 -\frac{d^2 \psi}{d r_*^2} +V(r) \psi = \omega^2  \psi
\end{equation}
Very loosely speaking, stability is attained if (\ref{RW3b}) has no solutions with {\em negative}
 $\omega^2$. More rigorously, this statement must be supplemented with the imposition of appropriate
 boundary conditions on $\psi(r)$. The resulting boundary value problem for $M > 0$ turns out to be
  quite different from that for $M<0$.

The case $M >0$ has been considered in detail in the literature
(see, e.g. \cite{chandra}). In this case we have $-\infty < r_* <
+\infty$ and $V(r)$ is positive definite in this range. Then,
 imposing boundary conditions that imply that the perturbed space time remains asymptotically flat,
  with a regular horizon, it can be shown that (\ref{RW3b})is turned into a self adjoint boundary
  value problem, that admits only solutions with $\omega^2 >0$, and, therefore, the spacetime is
  linearly stable.

For $M<0$ the problem is quite different. First, we have $0 \leq r_* < +\infty$.
 Second, as can be seen from (\ref{poten1}), $V(r)$ is singular for $r=r_*=0$,
 and for $\lambda r = -3 M$. Finally, since $r=0$ corresponds to a singularity,
 it is not at all clear what boundary conditions are physically adequate
 for $r_*=0$. An attempt at an appropriate formulation of the problem was recently
 given in \cite{ghi}. Unfortunately, for reasons explained below in Section \ref{ns}, we consider
 that there are some serious difficulties in that analysis that make their conclusions invalid.

\section{A modified Regge-Wheeler-Zerilli formalism}
\label{rwzg}

In view of the difficulties indicated in the previous Section, we present here a modification
of the Regge-Wheeler-Zerilli formalism that, as we shall show, allows for a simple analysis of
 the stability problem for $M<0$, with the intriguing result that these solutions are unstable
 under arbitrarily small perturbations.  Following \cite{ghi}, we define $\mu = -2M$, but consider,
 instead of $\psi$, the function,
\begin{equation}
\label{AppB1}
\zeta(t,r)= (2\lambda r -3 \mu ) \psi(r,t).
\end{equation}
which is also simply related to $H_0$, $H_1$, $H_2$, and $K$, and shares its uniqueness and gauge
 invariance properties.

 { } From (\ref{RW3}) we easily find that $\zeta$ satisfies the equation,
 \begin{eqnarray}
\label{RW17}
 {\frac {\partial ^{2} \zeta}{\partial {t}^{2}}} & = &
{\frac {   \left( r+\mu \right) ^{2}}{{r}^{2}}}\left( {\frac
{\partial^{2} \zeta}{\partial {r}^{2}}}  \right)   \nonumber \\ & &
-{\frac {\left( 4\lambda {r}^{2}+ 6 \lambda \mu r-3{\mu}^{2} \right)
\left( r+\mu
 \right)  }{{r}^
{3} \left( 2 \lambda r-3\mu \right) }}{\frac {\partial
\zeta}{\partial r}} -{\frac {
 \left( r+\mu \right)  \left( 3{\mu}^{2}- 6 \lambda \mu r+4 \lambda^2 {r}^{2}
 \right) }{{r}^{4} \left( 2 \lambda \,r-3\,\mu \right) }}\zeta
\end{eqnarray}
Since the coefficients in (\ref{RW17}) depend only on $r$, we may separate variables by
setting $\zeta(t,r)=\exp(i \omega t) \chi(r)$, which gives
\begin{eqnarray}
\label{AppB2}
{\frac {d^2 \chi}{d{r}^{2}}} & = &
 {\frac { \left(
4 \lambda {r}^{2}+6  \lambda r\mu -3{\mu}^{2} \right)  }{
 \left( 2  \lambda  r-3\mu \right)  \left( r+\mu
 \right) r}}{\frac {d \chi}{dr}} \nonumber \\ & &
+\left[{\frac { \left( 4{\lambda}^{2}{r}^{2}-6 \lambda
r\mu+3{\mu}^{2} \right)
  }{ \left( 2  \lambda  r-3\mu \right)
 \left( r+\mu \right) {r}^{2}}}  -\frac{r^2}{(r+\mu)^2}\omega^2 \right]{\chi}
\end{eqnarray}
This ODE, obviously, does not have the Schr\"odinger like form of
(\ref{RW3b}). Moreover, the coefficients of $\chi$ and $d \chi/dr$
are singular at $r=3 \mu /  (2 \lambda)$. However, the nature of
this singularity is rather subtle, and, as shown in detail in the
Appendix, it turns out that {\em all} solutions of (\ref{AppB2}) are
regular around $2 \lambda r=3\mu $, {\em irrespective} of the
boundary conditions for $r=0$ and $r=\infty$. In the next subsection
 we consider the {\em zero modes} or time independent solutions, as an illustration of this
 rather peculiar feature of (\ref{AppB2}).


\section{Static solutions}


Static solutions correspond to  $\omega=0$. In this case (\ref{AppB2}) is
reduced to,
\begin{equation}
\label{AppB3}
{\frac {d^2 \chi}{d{r}^{2}}}
 -{\frac { \left(
4 \lambda {r}^{2}+6  \lambda r\mu -3{\mu}^{2} \right)  }{
 \left( 2  \lambda  r-3\mu \right)  \left( r+\mu
 \right) r}}{\frac {d \chi}{dr}} -{\frac { \left( 4{\lambda}^{2}{r}^{2}-6
\lambda  r\mu+3{\mu}^{2} \right)
  }{ \left( 2  \lambda  r-3\mu \right)
 \left( r+\mu \right) {r}^{2}}}{\chi} =0
\end{equation}
A straightforward computation shows that this equation has, in general, two independent solutions, one
regular and the other with logarithmic divergences for $r=0$, and
for $r=-\mu$.  The regular
solution is a polynomial of order $\ell+2$, given by,
\begin{equation}
\label{AppB4}
 Q(r)=\sum_{i=1}^{\ell +2} q_i r^i
\end{equation}
with
\begin{equation}
\label{AppB5}
  q_{{i}}={\frac { \left( i-2 \right) \Gamma  \left( {\ell}+i-1 \right)  \left[  \left( i-4 \right)
   {\ell}\, \left( {\ell}
+1 \right) +i-1 \right] } { 3\, \Gamma( {\ell}-i+3) {\mu}^{i -1}
\left[ \Gamma( i )  \right] ^{2}}}
\end{equation}
The general solution of (\ref{AppB3}) may
be written in the form,
\begin{equation}
\label{AppB6}
\chi(r) = C_1 Q(r) + C_2 Q(r) \int_{C_3}^r{\frac{y (2
\lambda y -3 \mu)}{(y+\mu)Q(y)^2} dy}
\end{equation}
where $C_1$, $C_2$ and $C_3$ are constants. Notice that this
expression contains effectively only two arbitrary constants,
because a change in $C_3$ amounts only to a change in $C_1$. We also remark that, since all solution
of (\ref{AppB3}) are regular for $r \geq 0$, we must have $Q(r) > 0$ for $r>0$, for otherwise a zero
 in $Q(r)$ would introduce singularities in $\chi$ through the integral term in (\ref{AppB6}), which must be absent in the general solution of (\ref{AppB3}).

\subsection{Boundary conditions}

Equations (\ref{AppB5},\ref{AppB6}) provide the general time
independent solutions for the perturbative zero mode, but we still need to specify boundary conditions.
The question is
then if these solutions, perhaps for some particular choices of
$C_1$ and $C_2$, are acceptable as perturbations of the negative mass
Schwarzschild metric. As a possible criterion, we explore the change
in a particular curvature invariant, which we choose
 to be the Kretschmann invariant ${\cal{K}}$, given by,
\begin{equation}
\label{RW9} {\cal{K}} = R_{abcd} R^{abcd}
\end{equation}
where $R_{abcd}$ is the Riemann curvature tensor. Replacing
(\ref{AppB5})-(\ref{AppB6}) in the appropriate expression for
${\cal{K}}$, linearized to
 first order in the perturbations, we find that, in the $r \simeq 0$
 limit
\begin{multline}
\label{Kcerca} {\cal{K}} = \frac{12 \m^2}{r^6} - \epsilon \; C_2 \;
\left[  \frac{36 \m^2}{r^7}  + {\cal O} \left( \frac{\ln(r) }{r^5}
\right) \right] Y_{(\ell,m)}
\\ - \epsilon \; C_1 \;\left[ \frac{\ell (\ell ^2-1) (\ell +1)
(\ell+2)(\ell ^2+\ell +4)(\ell ^2+\ell-1)}{ 9 \; \m^2 r^3} + {\cal
O} \left( \frac{1}{r^2} \right)\right] Y_{(\ell,m)}
\end{multline}
 We must set $C_2=0$ to ensure that the first order term
does not grow faster than the zero order term as $r \to 0^{+}$. On
the other hand, if we  compute ${\cal{K}}$ for large $r$,  we obtain
\begin{multline}
\label{Klejos} {\cal{K}} = \frac {12 \mu^2} {{r}^{6}} -  \; \epsilon
\; C_1 \; \left[ \left( \frac{2 \; \Gamma (2 \ell +1)}{\mu^{\ell} \;
\Gamma (\ell-1)^2} \right) r^{\ell-5} +  {\cal O} (r^{\ell-4})
\right] \; Y_{(\ell,m)} \\
+ \epsilon \; C_2 \; \left[ \frac{K_{\ell}}{ r^{\ell+7}} +  {\cal O}
\left( \frac{1}{r^{\ell+8}} \right) \right] \; Y_{(\ell,m)},
\end{multline}
where
\begin{equation}
K_{\ell} = \frac{3 \sqrt{\pi} \; \Gamma(\ell)  (3 \ell^2+5 \ell
+8)(\ell+1)(\ell+2) \;\m^{\ell +2}}{ 4^{\ell} \; \Gamma(\ell +1/2)
(\ell-1)^2},
\end{equation}
then $C_1$ must also be set to zero in order to keep the
perturbation smaller than the zero order term  at large $r$.
 We must, therefore, conclude that there is
no (non trivial) choice of $C_1$ and $C_2$ for which a static
 perturbation will remain small in the entire range of $r$.\\





\subsection{A comment on a particular solution} \label{ns}

The $\ell = 2$ perturbation was studied in detail \cite{ghi}, where
it was shown  that (\ref{AppB3}) has the property that an appropriate choice $q_o$ of
 $q \equiv C_1/C_2$ in the general solution (\ref{AppB6}), gives a particular solution
 $\chi_o$, that is regular in $0 < r < +\infty$, and such that near $ r = 3 \mu/4$,
 the solution admits the expansion,
\begin{equation}
\label{RW15}
 \chi_o(r) \simeq C_1 \left( \a \left( 4\,r-3\,\mu \right) ^{3}+ \b
  \left( 4\,r-3\,\mu \right) ^{4}+ \dots \right)
\end{equation}
Then $\psi_o= \chi_o/( 4\,r-3\,\mu)$ satisfies $\psi_o=d\psi_o/dr=0$ for $r=3 \mu/4$, and one may construct a
bounded static solution ($\w^2=0$) of  (\ref{RW3b})  by the prescription,
\begin{equation} \label{nosol}
\hat \psi_o (r) = \begin{cases} \psi_o(r) & , 0 \leq r < 4\m/3 \\
0 & , r > 4\m/3 \end{cases}
\end{equation}
Departures of $C_1/C_2$ from $q_o$ in one particular direction would
then introduce a negative $\w^2$ in the spectrum of (\ref{RW3b}),
and thus an instability \cite{ghi}. We remark, however, that, although the
matching (\ref{nosol}) provides a solution
  of (\ref{RW3b}) for all $r > 0$, it is not a
  solution of the linearized Einstein equations, with which we are concerned here.
  The reason is that, in accordance with
(\ref{RW3}), the linearized Einstein equations include up to fourth
order $r-$derivatives of $\psi$, and   the second derivative of
(\ref{nosol}) has a discontinuity at $r=3 \mu /4$.
 Thus, inserting $\hat
\psi_o(r)$ in
 (\ref{RW3}) gives a perturbation that fails to be
 a solution of the linearized vacuum equations, because of the presence of Dirac's deltas,
  associated to the discontinuities of the higher derivatives. In
  the notation of equations (2.18)-(2.19) in \cite{ghi}, the Einstein
  tensor corresponding to (\ref{nosol})  has the following nontrivial
  components:
 \begin{eqnarray}
 G_{rr} &=& \left( \frac{256 \; C_4}{63}\right)   \frac{Y_{(2,m)}(\theta,\phi)}{r^2} \;\delta ( \mu - 4r/3) \\
G_{r \theta} &=& -\left( \frac{128 \; C_4}{63}\right)
\frac{\partial Y_{(2,m)}(\theta,\phi)}{\partial \theta}
\frac{1}{r^2} \;\delta ( \mu - 4r/3) \\
G_{r \phi} &=& -\left( \frac{128 \; C_4}{63}\right) \frac{\partial
Y_{(2,m)}(\theta,\phi)}{\partial \phi}
\frac{1}{r^2} \;\delta ( \mu - 4r/3) \\
G_{\theta \theta} &=& \left( \frac{128 \; C_4}{27}\right)\;
Y_{(2,m)}(\theta,\phi) \left[ \delta ( \mu - 4r/3) - r \delta' ( \mu
- 4r/3) \right] \\
G_{\phi \phi} &=& \left( \frac{32 \; C_4}{27}\right)\;
Y_{(2,m)}(\theta,\phi) \; sin^2(\theta) \;\left[ 3 \delta ( \mu -
4r/3) - 4 r \delta' ( \mu - 4r/3) \right]
\end{eqnarray}

\subsection{A physical interpretation for the zero modes}

Even though there are problems in interpreting the zero modes given
by (\ref{AppB4})-(\ref{AppB6}) as
 perturbations of the negative mass Schwarzschild metric, they are static solutions
  of the linearized equations and, therefore, their physical interpretation is of interest.
  One way in which the problem may be analyzed is to consider the Schwarzschild metric in
   Weyl coordinates. The general form of the static axially symmetric vacuum metrics in Weyl
   coordinates is,
\begin{equation}
\label{wey1}
ds^2=-e^{2U} dt^2+ e^{2k-2U} d\rho^2+\rho^2 e^{-2U} d\phi^2 +e^{2k-2U} dz^2
\end{equation}
where $U=U(\rho,z)$, $k=k(\rho,z)$, and the Einstein equations imply,
\begin{equation}
\label{wey2}
\frac{\partial^2 U}{\partial \rho^2} +\frac{1}{r}
\frac{\partial U}{\partial \rho}+\frac{\partial^2 U}{\partial z^2}=0
\end{equation}
and,
\begin{eqnarray}
\label{wey3}
 k,_{\rho} & = & \rho \left(U,_{\rho}\right)^2-  \rho \left(U,_{z}\right)^2 \nonumber \\
k,_{z} & = & 2 \rho  U,_{\rho}U,_{z}
\end{eqnarray}
where a comma indicates partial derivative.
It can be checked that (\ref{wey2}) is the integrability condition for the system (\ref{wey3}),
so that $k$ is obtained by quadratures if $U$, satisfying (\ref{wey2}) is given. The latter is
the Laplace equation in flat cylindrical coordinates and, as is well known, allows for the
construction of static vacuum solutions of the Einstein equations by taking for $U$ different
 solutions of the Laplace equation corresponding to localized sources. The Schwarzschild space
 time is recovered by taking for $U$ the Newtonian potential corresponding to uniform line
  source of finite length, centered at the origin, and laying along the z-axis. The explicit
   forms for $U$ and $k$ in this case are,
\begin{eqnarray}
\label{wey4}
U_S(\rho,z) & = & \frac{1}{2} \ln\left[\frac{\rho_+ +\rho_- -2M}{\rho_+ +\rho_- +2M}\right] \nonumber \\
k_S(\rho,z)& = &  \frac{1}{2} \ln\left[\frac{(\rho_+ +\rho_-)^2-4M}{ 4 \rho_+ \rho_-}\right]
\end{eqnarray}
where $\rho_{\pm}=\sqrt{\rho^2+(z\pm M)^2}$. The constant $M$ is identified with the Schwarzschild
mass and can be positive or negative. Suppose now that instead of $U_S$ we consider, e.g.,
\begin{equation}
\label{wey5}
U(\rho,z) = U_S(\rho,z)+U_1(\rho,z;m)
\end{equation}
where $U_1(\rho,z;m)$ is a solution of (\ref{wey2}) that depends smoothly on a real parameter
$m$ such that,
\begin{equation}
\label{wey6}
 U_1(\rho,z;0)= 0
\end{equation}
This implies that $U(\rho,z)$ given (\ref{wey5}) satisfies also (\ref{wey2}) and can be used
to construct static vacuum solutions of the Einstein equations. But, in accordance with
(\ref{wey6}), these solutions, at least locally, can be made to approach the Schwarzschild
metric arbitrarily closely, by taking $m$ small enough. This may fail, and in general it does,
 when we approach the line source of $U_S$. This line source corresponds to the regular horizon
 for $M>0$ and, therefore, a singularity in $U_1$ destroys that regularity, and $U_1$ cannot be
  considered globally as a ``small" perturbation. Similarly, for $M<0$, the line source
  corresponds to a naked singularity, and a singularity in $U_1$ may modify its structure,
  and can make it even ``harder''. To illustrate this discusion we consider for $U_1$ the form,
\begin{equation}
\label{wey7}
U_1(\rho,z;m) = \frac{m}{\sqrt{\rho^2+z^2}}
\end{equation}
This represents (in Weyl coordinates) a point source at the origin. The resulting equations
 for the modified $k(\rho,z)$ may be solved explicitly, for any $m$, but we are only interested
 in the change to order $m$. Actually what we are interested in is the form of the metric
 in Schwarzschild, rather than Weyl coordinates. The coordinate change that accomplishes
 this is well known (see, e.g. \cite{exactas}, for details). Performing it we find that
 in Schwarzschild coordinates the modified metric, to first order in $m$, takes the form,
\begin{eqnarray}
\label{wey8}
g_{tt} & = & -1+ {\frac {2M}{r}}-{\frac {2 m \left( r-2M \right) }{r
\sqrt {{M}^{2} \cos^2 \theta   +{r}^{2}
-2Mr}}}
\nonumber \\
g_{rr} & = & {\frac {r}{r-2 M}}+{\frac {2 mr \left( 2r-3M-2
\sqrt { {M}^{2} \cos^2 \theta +{r}^{2}
-2Mr} \right) }{ \left(r -2M\right) \sqrt { {M}^{2}\cos^2
\theta    +{r}^{2}-2Mr}M}}
\nonumber \\
g_{\theta \theta} & = & {r}^{2}+{\frac {2 m {r}^{2} \left(2 r -3M-2\sqrt {
 {M}^{2} \cos^2 \theta   +{r}^{2}-2Mr}
 \right) }{M\sqrt {M^2  \cos^2  \theta +{r}^{2}-2Mr}}}
\nonumber \\
g_{\phi \phi} & = &  r^2 \sin^2  \theta -{\frac {2 m r^2
  \sin^2  \theta  }{
\sqrt { M^2 \cos^2 \theta +{r}^{2}
-2Mr}}}
\end{eqnarray}
This metric may be expanded now in the Regge-Wheeler-Zerilli functions. For the metric
in the form (\ref{wey8}), in the notation of \cite{RW}, the only non vanishing functions
are $H_0$, $H_2$, $K$, and $G$, corresponding to even (scalar) perturbations for even $\ell$, as the result of the  symmetry of $U_1(\rho,z;m)$ under $ z \leftrightarrow -z$.
This is not in the Regge-Wheeler gauge, but, from gauge invariance \cite{moncrief},
the function $\psi(r,t)$ is in this case given by,
\begin{eqnarray}
\label{wey9}
\psi(r,t) & = & \frac{1}{\lambda r+3M} \chi(r,t) \nonumber \\
\chi(r,t) & = & \frac{2r \left(r -2M\right)}{\ell (\ell+1) }  \left( H_2   -r{
\frac {\partial K}{\partial r}}   -{\frac { \left( r-
3M \right) K  }{r-2M}} \right)  + {r}^{2}  \left( K   +
 \left(r -2M \right) {\frac {\partial G}{\partial r}}    \right)
\end{eqnarray}
Going back to (\ref{wey8}), if we specialize to $\ell=2$, we find,
\begin{eqnarray}
\label{wey10}
H_2 & = & {\frac {\sqrt {
5 \pi } m \left( 2r-3M\right)  }{2{M}^{4}}} \left[  \left( 2\,{M}^{2}-6
\,Mr+3\,{r}^{2} \right)  \ln  \left(1 -\frac{2M }{r}\right) +6\, \left( r-M \right) M \right]
  \nonumber \\
K & = &  {\frac {\sqrt {5 \pi } m }{4{M}^{4}}}\left[ \left( 3\,{r}^{
3}-15\,{r}^{2}M+16\,r{M}^{2}-2\,{M}^{3} \right)   \ln  \left(1 - \frac{2M}{r}
\right) -24\,r{M}^{2}+10\,{M}^{
3}+6\,{r}^{2}M \right]
\nonumber \\
G & = & {\frac {\sqrt {5 \pi }m}{12{M}^{4}}} \left[ 3 \left(M- r
 \right)  \left( 2{M}^{2}-2Mr+{r}^{2} \right)  \ln
 \left(1 -\frac{2M}{r} \right)
+2 M \left(6 r M -7 M^2 -3 r^2 \right)
  \right]
\end{eqnarray}
and the corresponding expression for $\chi(r)$ results,
\begin{equation}
\label{wey11}
\chi= {\frac {  \sqrt {5 \pi } m r }{6{\mu}^{3}}}\left[ 3 \left( 4{r}^{3}-6{r}^{2}\mu+3{\mu}^{3}
 \right)   \ln  \left( 1+\frac{\mu}{r}
 \right)  -12\mu {r}^{2}+24\mu^2 r-13{\mu}^{3}
 \right]
\end{equation}
which coincides with (\ref{AppB4})-(\ref{AppB6}) for  an appropriate
choice of integration constants. Similar results hold for the other
even values of $\ell$. Non vanishing contributions for odd $\ell$ may be
 obtained starting from a different choice of $U_1(\rho,z;m)$ that
 makes it odd under $z \longleftrightarrow -z$, with analogous results
 regarding the corresponding zero modes. Therefore, the zero modes are the linearized
expressions for the modifications that result in the Schwarzschild
metric by the addition of different (static) multipoles. We stress
again that these modifications approach only locally the
Schwarzschild metric, because they modify non trivially the structure of its
singularity. We shall add more comments on this feature in Section
\ref{cc}.


\section{Time dependent solutions}

Going back to (\ref{AppB2}), we consider now the general, time dependent case, with
 $\omega \neq 0$. We are interested in solutions that are regular in $0\leq r \leq \infty$.
 Therefore, we need to consider the behavior of $\chi_1$ near the singular points
  $r=0$, $r=3\mu/4$ and $r=\infty$. First we find that near $r=0$,
  any solution of (\ref{AppB2}), for any $\omega$, is of the form,
\begin{equation}
\label{RW20}
\chi_1(r) \simeq C_1 r + C_2 r \ln(r) + \dots
\end{equation}
where $C_1$ and $C_2$ are constants, and $\dots$ indicates terms that vanish faster
for $r=0$. This implies that all solutions of (\ref{AppB2}) vanish for $r=0$, but,
in general, $d \chi_1/dr$ diverges as $\ln(r)$. On account of our previous discussion
regarding acceptability of the perturbations, this means that we need to impose the condition
that this divergence is absent, and therefore, only the regular solution associated with $C_1$
can be considered for the stability analysis. This fixes our boundary condition for $r=0$.

Next, we have to consider the singularity for $r=3\mu/4$. Here, as already mentioned,
and as shown in Appendix A, we have the rather surprising feature that {\em all}
 solutions of (\ref{AppB2}) are analytic for $r=3\mu/4$, and, therefore, we need only
 be concerned with the behavior of $\chi_1$ for large $r$, as no restriction is imposed
 by the behavior at $r=3\mu/4$.

The behavior at $r=\infty$ is also simple. For large $r$, the general solution of (\ref{AppB2})
 is of the form
\begin{equation}
\label{RW21}
 \chi_1(r) \simeq   \left( {C_3} r +\dots
 \right) \exp(i r\sqrt{\omega^2})+  \left( C_4 r +\dots\right)
\exp(-i r\sqrt {\omega^2})
\end{equation}
where $C_3$ and $C_4$ are constants, and, again, $(\dots)$ indicate
sub leading terms \footnote{ Notice that the leading term in $r$ in
(\ref{RW21}) can be traced to the
 transformation (\ref{AppB1}), since, for any sign of $M$, the Zerilli equation for $\psi(t,r)$
 approaches the one dimensional wave equation, and $\psi(t,r) \rightarrow \exp(i \omega t)$.}

From the assumed time dependence,
positive values for $\omega^2$ correspond to an oscillating, and
therefore stable, behavior for the perturbations. The interesting
range for instability is, therefore, $\omega^2 < 0$.
 On account of (\ref{RW21}), the solutions for $\omega^2 < 0$ will
grow unboundedly for large $r$, unless we can eliminate the exponentially divergent terms.
  In the following section we analyze this problem in detail.

\section{Unstable solutions} \label{us}

Unstable modes, correspond to solutions of (\ref{AppB2}) with $\omega^2 =-k^2$, that
 satisfy the regularity condition for $r=0$, and are exponentially decreasing for large $r$.
  Therefore, finding these solutions  turns this equation into a boundary value problem,
  determining the appropriate value or values of $k$, if they exist.

As a first step, we considered a ``shooting'' approach, and, using a high precision
numerical method, we found, for $\ell=2$, a numerical solution for $k \simeq 4/\mu$.
Since this value was so close to an integer times $1/\mu$, we
replaced $k = 4/\mu$, $\ell=2$ in (\ref{AppB2}), and found that indeed it has a rather
simple solution for this value of $k$. Repeating the procedure for $\ell=3$, we found again
a solution for $k=20/\mu$. This suggested the possibility that we could construct solutions
 for {\em all} $\ell$, by a simple generalization of the results for $\ell=2$, and $\ell=3$.
 This was effectively the case, and we found that,
\begin{equation}
\label{AppC3}
 {\chi} ( r )  = r{e^{-kr}} \left( r+\mu \right) ^{k\mu}
 \left[ {C_1}+{C_2}\,\int _{{C_3}}^{r}\!{\frac { \left( 2\,
\lambda\,y-3\,\mu \right) ^{2}{e^{2\,ky}}}{ \left( y+\mu \right) ^{2\,
k\mu+1}y}}{dy} \right]
\end{equation}
where $C_1$, $C_2$, and $C_3$ are arbitrary constants, is a solution of (\ref{AppB2}) with
$\omega^2=-k^2$, provided  the constant $k$ is given by
\begin{equation}
\label{AppC4}
 k = \frac{2\lambda (\lambda+1)}{3 \mu}=\frac{ (l-1) l (l+1) (l+2)} {6 \mu}
\end{equation}
and, therefore, $k\mu$ is an integer for all $\ell$. Notice that (\ref{AppC3}) effectively
 contains only two arbitrary constants, because any change in $C_3$ can be compensated by a
 corresponding change in $C_1$.

It is easy to see that the integral in (\ref{AppC3}) contributes logarithmic divergences for
 $r=0$, and an exponentially increasing factor for $r \rightarrow \infty$. Therefore if we
 choose $C_2=0$, we find that
\begin{equation}
\label{AppC5}
 {\zeta} ( t, r )  = C_1 r{e^{-kr}} \left( r+\mu \right) ^{k\mu}e^{kt}
\end{equation}
with $k$ given by (\ref{AppC4}), is a solution of (\ref{AppB2}), representing
 the general expression for the regular, bounded, unstable modes.
Moreover, this result implies that {\em all scalar modes (i.e. for all $\lambda$) of the Schwarzschild negative
mass naked singularity are linearly unstable}.

As a check of the consistency of our results, we replace (\ref{AppC5}) in (\ref{AppB1}),
 and this in(\ref{RW2}), and obtain
\begin{eqnarray}
\label{RW24a}
 K(t,r) & = & -\frac{ (\lambda + 1)(r+\mu)^{k \mu}}{3 \mu}\exp[k(t- r)]   \nonumber \\
 H_1(r,t) & = & -H_2(t,r) = -\frac{\lambda (\lambda+1)
   [2(\lambda+1) r+ 3 \mu] \; r  (r+\mu)^{k\mu -1}}{9 \mu^2}\exp[k(t- r)] \nonumber
\end{eqnarray}
which shows explicitly the regularity of the solution of the perturbative equations in the
 range $0 \leq r < +\infty$. The perturbations of the curvature are also well behaved.
Replacing in the Kretschmann scalar, we find that it is unchanged from the unperturbed value,
to order $\epsilon$.
Furthermore, a complete set of quadratic algebraic invariants of the
Riemann tensor \cite{cm,zm} was computed to find  that they are all
invariant to first order in $\epsilon$. However, we have found non
trivial first order contributions to the  {\em differential}
invariants
\begin{eqnarray} \label{i1}
R_{abcd;e} \; R^{abcd;e} &=& \frac{180 \mu^2 (r + \mu)}{r^9} - \e \;
  \frac{ 5 \ell (\ell-1) (\ell+1)(\ell+2)}{r^8}  \;  Y_{\ell m}(\theta,\phi)\\
\nonumber && \times  \left[ r \ell (\ell+1) + 3 \mu \right]
  (r+\mu)^{k \mu} \;
 \exp[k (t-r)]
\end{eqnarray}
and \begin{multline}
 \label{i2}
R_{abcd;ef} \; R^{abcd;ef} = \frac{90 \m^2 (65 \m^2 +120 \m r+56
r^2)}{r^{12}} \\- 30 \e k \frac{(r+\mu)^{k \mu} \exp[k
(t-r)]}{r^{11}} \; Q_{\ell}(r) \;Y_{\ell m}(\theta,\phi)
\end{multline}
\begin{multline}
Q_{\ell}(r) = \left( {l}^{2} \left( l-1 \right)  \left( l+2 \right)
\left( l+1 \right) ^{2 }{r}^{3}+3\,\mu\,l \left( l+1 \right)  \left(
{l}^{2}+l+14 \right) {r} ^{2} \right. \\ \left.+ 3 \,{\mu}^{2}
\left( 17\,{l}^{2}+17\,l+56 \right) r+180\,{\mu}^{3} \right)
\end{multline}
and we may expect non trivial contributions in some of the many other differential
invariants. Note that, both for  (\ref{i1}) and (\ref{i2}),
the ${\cal O} (\e)$ term  grows slower than the ${\cal O} (\e^0)$ term as
$r \to 0^+$ limit, and also as $r \to \infty$.

\section{Comments and conclusions} \label{cc}

In this article we have shown that even if we only allow for the
 most restrictive boundary conditions, both as $r \to 0^+ $ and $r \to \infty$, the
 perturbation equations for the negative mass Schwarzschild metric contain,
 for every angular mode, solutions that grow exponentially in time.
 These solutions represent everywhere regular perturbations that can be
 made initially globally arbitrarily small and, therefore, imply that
 the metric is unstable at the perturbative level, for all angular modes.
 An interesting technical point of our derivations was the introduction of
a modification of the standard Regge-Wheeler-Zerilli method, that led to equation
 (\ref{AppB2}) and its rather unusual properties.
Our analysis of the zero modes indicates that rather than
perturbations they should be considered as linearized approximations
for exact static metrics that locally, but not globally approach the
negative mass Schwarzschild metric. Although asymptotically flat,
they actually contain a naked singularity where the curvature
diverges even faster than for the Schwarzschild metric. From the
perspective of the present analysis, these spacetimes should also be considered
unstable, since, if they are considered as part of the unperturbed
metric, one would be allowed to relax the boundary conditions at
$r=0$, and as can be seen from the derivations in Section \ref{us}, this
would lead to unstable solutions for {\em any} $k = \sqrt{-\omega^2} $. Our
results also suggest that, even though the vector (linearized)
perturbations have been shown to be stable \cite{ghi}, one would
expect that these would also become unstable through the coupling to
the scalar modes. Since vector modes are related to rotation, this
again suggests that rotation would not make the singularity more
stable, although going through this in detail would require higher
order perturbation theory, or, perhaps, some application of the
Teukolsky equation.

Finally, there is an obvious implication of our results to the
cosmic censorship conjecture, since the linear instability implies
that the negative mass \s cannot be the final stage of
an evolving spacetime.

\appendix

\section{Behavior of the ODE (\ref{AppB2}) around $r = (3 \mu)/(2 \lambda)$ }

In this Appendix we give a brief analysis of the behavior of
equation (\ref{AppB2}) near $r~=~3\mu/(2 \lambda)$. For this purpose
it is simpler to change variables to $\rho=r-3\mu/(2 \lambda)$, so
that the singularity is located at $\rho=0$. Replacing $\rho-$
derivatives for $r-$ derivatives, (\ref{AppB2}) takes the form
\begin{equation}\label{ed}
 \frac {d^2\chi}{d\rho^2} +  \frac {Q(\rho)}{\rho}\frac {d
\chi}{d\rho} +\frac {P(\rho)}{\rho}\chi = 0,
\end{equation}
where $Q$ and $P$ admit expansions
\begin{equation}
Q(\rho) = \sum_{k\geq 1} a_k \rho^{k-1} \hspace{2cm} P(\rho) =
\sum_{k\geq 1} b_k \rho^{k-1}
\end{equation}
with $a_1 \neq 0$ and $b_1 \neq 0$. Equation (\ref{ed}) has, in general,
two linearly independent solutions. Let us assume there is a solution
which is analytic around $\rho=0$:
\begin{equation}
\chi(\rho) = \sum_{k\geq 0} c_k \rho^{k},
\end{equation}
then (\ref{ed}) reduces to
\begin{equation} \label{super}
s \left( s-1 + a_1\right) c_{{s}}+\sum _{l=0}^{s-1} \; c_{{l}} \;
\left( l \; a_{{s+1-l}} +  b_{{s-l}} \right)  = 0 , \;\; s=1,2,3,...
\end{equation}
There are different possibilities:\\
{\bf \em (A) $a_1$ is not a negative integer:} in this case
(\ref{super}) can be solved iteratively (solve the $s=1$ equation,
then the $s=2$ one, etc) starting from an arbitrary $c_0$ value.
This gives one solution
of (\ref{ed}) which is analytic around $\rho=0$.\\
{\bf \em (B) $a_1$ is  a negative integer:} let $a_1=1-s_o$. The $s
\leq s_o$ equations (\ref{super}) constitute a system of $s_o$
homogeneous linear equations on $s_o$ unknowns ($c_0$ through
$c_{s_o-1}$). Then, we have two separate cases depending on whether
the determinant of the system vanishes or not.

  {\bf (case {\em B.1})} If the determinant of
this system is nonzero, the only solution
 is $c_0=c_1=\dots=c_{s_o-1}=0$. The $s
> s_o$ equations are then solved iteratively starting from an
arbitrary $c_{s_o}$ value. Again, this gives one solution of
(\ref{ed}) which is analytic around $\rho=0$ (moreover, $\chi =
{\cal O} (\rho^{s_o})$).

{\bf case {\em B.2}} If the
above mentioned determinant is zero, the $s_o$ equation will be
consistent with the previous ones, and will add no information.
$c_{s_o}$ can be given an arbitrary value and the remaining ($s
> s_o$) equations solved iteratively in terms of the {\em two}
arbitrary parameters $c_0$ and $c_{s_o}$.\\

Thus, generically, in case $B.1$, (\ref{ed}) admits only one solution analytic
around $\rho=0$. However, in the special case $B.2$ {\em both}
solutions of (\ref{ed}) are analytic around $\rho=0$. \\
If we expand
(\ref{AppB2}) we get
   \begin{multline} \label{exp}
0 = \frac{d^2 \chi}{d \rho^2} + \left( -2\,{\rho}^{-1}-{\frac
{4{\lambda}^{2}}{3 \mu\, \left( 3+2\,\lambda
 \right) }}+\,{\frac {16 \left( 3+\lambda \right) {
\lambda}^{3}\rho}{9{\mu}^{2} \left( 3+2 \lambda
 \right)^2 }} + {\cal O} (\rho^2) \right) \frac{d \chi}{d \rho} \\+
 \left({\frac {-4{\lambda}^{2}}{3 \mu\, \left( 3+2\,\lambda \right) \rho}}-
\,{\frac {16\,{\lambda}^{4}-81\,{w}^{2}{\mu}^{2}}{9{\mu}^{2} \left(
3+2\lambda \right)^2 }}+\,{\frac {8{ \lambda}^{2} \left(
4\,{\lambda}^{3}+27\,{w}^{2}{\mu}^{2} \right) \rho }{9{\mu}^{3}
  \left( 3+2\lambda \right)^3 }} + {\cal{O}}(\rho^2) \right)
\chi(\rho)
 \end{multline}
 Since $a_1=-2$,  this equation belongs to case $B$ with $s_o=3$.
 In order to know whether this is a case $B.1$ or a case $B.2$,
 we need to calculate the determinant of the matrix associated to
 the system of equations (\ref{super}) with $s=1,2,3$, i.e.
 \begin{equation} M = \left(
 \begin{matrix} b_1 & a_1 & 0 \\ b_2 & b_1+a_2 & 2 a_1+2\\
 b_3 & a_3+b_2 & b_1+2 a_2 \end{matrix}  \right)
 \end{equation}
Inserting the $a_i$ and $b_i$ values from (\ref{exp}), we find that
the matrix $M$ above has zero determinant, {\em for any value of
$\lambda$ and  $\w$}. We conclude that our system belongs to case
$B.2$, then both solutions of (\ref{AppB2}) are analytic, and this
holds for any
$\w$.\\

An example of a differential equation (\ref{ed}) of type A is
\begin{equation}
\frac {d^2\chi}{d\rho^2} -  \left(\frac {3}{2 \rho} \right)\frac {d
\chi}{d\rho} +\frac {1}{\rho}\chi = 0.
\end{equation}
This has solution
\begin{multline}
\chi(\rho) = C_1 \left[ 6 \sqrt{\rho}\; \sin(2 \sqrt{\rho} ) + (3-4
\rho) \cos(2 \sqrt{\rho} ) \right]  + C_2 \left[ 6 \sqrt{\rho}\;
\cos(2 \sqrt{\rho} ) + (4 \rho -3) \sin(2 \sqrt{\rho} ) \right] \\ =
C_1 \left[ 3 + 2 \rho + 2 \rho^2 -\frac{4}{3} \rho^3 + \frac{2}{9}
\rho^4 + \cdots  \right] + C_2 \left[ -\frac{32}{15} \rho^{5/2} +
\frac{64}{105} \rho^{7/2} - \frac{64}{945} \rho^{9/2} + \cdots
\right]
\end{multline}
An example of a differential equation (\ref{ed}) of type B.1 is
\begin{equation}
\label{app1a} \frac {d^2\chi}{d\rho^2} -\frac{2}  {\rho}\frac {d
\chi}{d\rho} +\frac {b_1}{\rho}\chi = 0
\end{equation}
This has the general solution
\begin{equation}
\label{app1b} \chi  ={C_1}\,{\rho}^{3/2}{J_3}
 \left( 2\,\sqrt { b_1 \rho} \right) +{C_2}\,{\rho}^{3/2}{
Y_3} \left( 2\,\sqrt { b_1 \rho} \right)
\end{equation}
where $J_i$ and $Y_i$ are Bessel functions. One solution is
analytic, whereas the other contains  terms
 of the form $\rho^3 \ln(\rho)$, so that it is only twice differentiable at $\rho=0$.

Finally, case $B.2$ is illustrated by the equation,
\begin{equation}
\label{app6}
\frac {d^2\chi}{d\rho^2}
+\left(-\frac{2}  {\rho}+a_2\right)\frac {d \chi}{d\rho}
+\left(\frac {a_2}{\rho}- \frac{3 a_2^2}{4}\right)\chi = 0
\end{equation}
which has the general solution,
\begin{equation}
\label{app6a}
 \chi(\rho) ={ C_1}\,\exp(a_2 \rho/2)+{
C_2}  \left( 1+2 a_2\rho+2 a_2^{
2}{\rho}^{2} \right) \exp(-a_2 \rho/3)
\end{equation}
which is clearly analytic for $\rho=0$.


\section*{Acknowledgments}
 G.D. thanks the KITP at UCSB, where part of this work was done,
 for  hospitality.
This work was supported in part by grants of CONICET (Argentina) and
the Universidad Nacional de   C\'ordoba. It was also supported in
part by grant NSF-INT-0204937 of the National Science Foundation of
the US. The authors are supported by CONICET (Argentina).




\begin{thebibliography}{99}
\bibitem{pm} Carroll S M (2003) {\em Phys. Rev.} {\bf C68} 023509,
astro-ph/0301273.\\
Gibbons  G W (2003), hep-th/0302199.\\
Caldwell R R (2002) {\em Phys. Lett.} {\bf B545} 23,
astro-ph/9908168.
\bibitem{ghi} Gibbons G W, Hartnoll D and  Ishibashi A 2005 {\em Prog.Theor.Phys.}
{\bf 113} 963-978, hep-th/0409307.
\bibitem{RW} Regge T and  Wheeler J A 1957 {\it Phys. Rev.} {\bf 108}, 1063.
\bibitem{Zerilli}  Zerilli F J 1970 {\it
Phys. Rev. Lett} {\bf 24} 737.
\bibitem{chandra} Chandasekhar S, {\em The mathematical theory of
black holes}, Oxford University Press, 1992.
\bibitem{cm} Carminati J  and McLenaghan R G 1991  {\it Jour.Math.Phys} {\bf 33},
3135.
\bibitem{zm} Zakhary E and  McIntosh C B G 1997 {\em Gen.Rel.Grav} {\bf 29}, 539.
\bibitem{exactas} Stephani H, Kramer D, MacCallum M, Hoenselaers C and Herlt E {\em
Exact Solutions to Einstein's field equations} 2nd ed (2003).
\bibitem{moncrief} Moncrief V (1974) {\em Ann. Phys} {\bf 88} 323.
\end{thebibliography}


\end{document}

The expression: "in physics, what is not forbidden by a physical law must necessarily occur'', or something similar to this has been sometimes attributed to R. Feynmann, and ...
